% ============================================================
% Chapter 3: [Your Chapter Title]
%
% PURPOSE:
%   This is the first of your "core" chapters — the ones that
%   present your actual contribution. What goes here depends
%   entirely on your thesis type. Choose the structure that
%   fits your work and rename the \chapter{} title below.
%
% COMMON OPTIONS FOR THIS CHAPTER:
%
%   "Methodology"
%     Describe how you conducted your research: your design,
%     data sources, tools, and evaluation plan. Use this for
%     experimental, ML, or data-driven theses.
%
%   "System Design"
%     Describe the architecture, components, and design decisions
%     of a system you built. Use this for software engineering
%     or implementation-heavy theses.
%
%   "Analysis" or "Synthesis"
%     Present a structured analysis or synthesis of findings
%     drawn from your literature review. Use this for research
%     or survey theses that do not involve building a system.
%
%   "Framework" or "Proposed Approach"
%     Define a new conceptual framework, model, or approach.
%     Use this for theses that propose a new way of thinking
%     about a problem.
%
% TIPS:
%   - Be specific enough that another student could reproduce
%     your work based on this chapter alone.
%   - Justify your choices: explain not just what you did
%     but why you did it that way.
%   - Use figures to illustrate architectures or workflows.
%     See LATEX_GUIDE.tex for figure examples.
%   - Use equations where appropriate.
%     See LATEX_GUIDE.tex for equation examples.
%
% TARGET LENGTH: 8--15 pages.
% ============================================================

\chapter{[Your Chapter Title Here]}
\label{chapter3}

\noindent
[Write a brief overview of this chapter. Describe what you present
here and how it connects to the research questions from Chapter 1.]


\section{[First Section]}

\noindent
[Write your first section here. If this is a methodology chapter,
describe your overall research design and justify your approach.
If this is a design chapter, describe your system at a high level
before breaking it down into components.]

% A figure showing an overview, architecture, or workflow is
% strongly recommended in most thesis types. Example:
%
% \begin{figure}[H]
%     \centering
%     \includegraphics[width=0.8\textwidth]{Images/your_figure.png}
%     \caption{A descriptive caption explaining what the figure shows.}
%     \label{fig:overview}
% \end{figure}
%
% Reference it in text like this: Figure~\ref{fig:overview} shows...


\section{[Second Section]}

\noindent
[Write your second section here. If this is a methodology chapter,
describe your data sources, participants, or experimental setup.
If this is a design chapter, describe individual components.]


\section{[Third Section]}
% Add or remove sections as needed.

\noindent
[Write your third section here. If this is a methodology chapter,
describe your evaluation metrics and how you will measure success.
If this is a design chapter, describe implementation considerations.]

% If you have an equation, use the equation environment:
%
% \begin{equation}
%     \text{Accuracy} = \frac{TP + TN}{TP + TN + FP + FN}
%     \label{eq:accuracy}
% \end{equation}
%
% Reference it in text like this: Equation~\ref{eq:accuracy} defines...
